\chapter{Zinc}

\section*{What Is This Test?}
The zinc test measures the concentration of zinc in serum or plasma. Zinc is an essential trace element that serves as a catalytic or structural cofactor for more than 300 enzymes and 2,000 transcription factors. It is required for immune function (T-cell development and function), wound healing, protein and DNA synthesis, cell division, taste and smell perception, and the function of superoxide dismutase (antioxidant defense). Zinc is absorbed in the proximal small intestine, transported in blood bound primarily to albumin (approximately 70\%) and alpha-2-macroglobulin, and excreted mainly through the gastrointestinal tract, with only 1--2\% lost in urine. Because the body has no dedicated zinc storage pool and relies on dietary intake and regulated absorption, serum zinc is a useful but imperfect marker of zinc status. Approximately 17\% of the world's population is at risk of zinc deficiency, most commonly due to inadequate dietary intake, malabsorption, chronic diarrhea, or increased losses. Serum zinc is the most widely used laboratory measure of zinc status, though it does not always correlate with whole-body zinc stores because of homeostatic redistribution during acute illness and inflammation.

\section*{Alternative Names}
Serum Zinc; Plasma Zinc; Zinc Level; Zn; Trace Element Panel (component)

\section*{Principle}
Serum or plasma zinc is measured by inductively coupled plasma mass spectrometry (ICP-MS), the current reference method, or by flame atomic absorption spectroscopy (AAS) at 213.9 nm. In ICP-MS, the sample is nebulized and ionized in an argon plasma, and zinc isotopes (Zn-66, Zn-68) are quantified by mass-to-charge ratio; collision-cell technology is used to resolve interferences. In flame AAS, the sample is aspirated into a flame and the absorbance of zinc at 213.9 nm is measured against standards. Pre-analytical contamination is the dominant source of error: samples must be collected in trace-element-free (royal blue-top) tubes with certified low zinc content, because rubber stoppers and glass tubes can leach zinc. Hemolyzed samples must be rejected because red blood cells contain approximately ten times more zinc than serum, and released zinc falsely elevates the result. Capillary (fingerstick) specimens are avoided because contact with the skin surface contaminates the sample.

\section*{History}
Zinc was recognized as essential for human health in 1961 when Ananda Prasad described the syndrome of growth retardation, hypogonadism, and delayed maturation in zinc-deficient Iranian and Egyptian adolescent males. The development of atomic absorption spectroscopy in the 1960s provided the first reliable method for measuring serum zinc, and flame AAS became the clinical standard for decades. The recognition that low zinc contributes to impaired wound healing, immune dysfunction, and diarrhea in developing countries drove population-level interest in zinc supplementation. Inductively coupled plasma mass spectrometry, developed commercially in the 1980s, later became the reference method because of its sensitivity and multi-element capability, allowing zinc to be measured alongside copper, selenium, and other trace elements in a single panel.

\section*{Why This Test Is Ordered}
\begin{itemize}
  \item Evaluation of suspected zinc deficiency in patients on prolonged parenteral nutrition, which classically produces acrodermatitis enteropathica-like skin lesions and is a medical emergency
  \item Assessment of patients with malabsorption syndromes: Crohn disease, celiac disease, cystic fibrosis, short bowel syndrome, and chronic pancreatitis
  \item Investigation of unexplained growth retardation, hypogonadism, or delayed sexual maturation in children and adolescents
  \item Evaluation of chronic diarrhea, which causes significant enteric zinc losses
  \item Assessment of patients with alcohol use disorder, cirrhosis, or chronic liver disease, who are frequently zinc deficient
  \item Evaluation of impaired wound healing or recurrent infections, especially in surgical and critically ill patients
  \item Investigation of unexplained hair loss, taste disturbances (hypogeusia), or skin changes
  \item Monitoring for zinc toxicity in patients taking high-dose zinc supplements or with occupational exposure
  \item Assessment of patients with acrodermatitis enteropathica, a rare autosomal recessive disorder of intestinal zinc absorption
  \item Co-measurement in Wilson disease evaluation, where serum copper and ceruloplasmin are low but zinc may be elevated
\end{itemize}

\section*{Primary vs Secondary Test}
Zinc measurement is a secondary test, ordered when deficiency or toxicity is clinically suspected rather than as a routine screening test. Because serum zinc falls during the acute-phase response and may not reflect tissue stores, it is best interpreted alongside clinical findings, dietary assessment, and, when indicated, functional measures.

\section*{How This Test Is Performed}
A venous blood sample is collected by standard venipuncture directly into a trace-element-free (royal blue-top) tube containing EDTA or lithium heparin. The tube must be filled completely and the sample processed promptly. Morning fasting samples are preferred because serum zinc varies modestly with meals (postprandial levels fall 10--20\%) and follows a circadian rhythm. The laboratory measures zinc using ICP-MS or atomic absorption spectroscopy. Whole blood, erythrocyte, and urine zinc can also be measured in specialized circumstances: erythrocyte zinc reflects longer-term status, and a 24-hour urine zinc is used to assess gastrointestinal losses or to monitor treatment in acrodermatitis enteropathica.

\section*{What Preparation Is Needed}
Fasting for 8--12 hours is preferred because food transiently lowers serum zinc. Collection must be performed with certified trace-element-free tubes and acid-washed collection equipment; routine tubes and some evacuated collection systems can contaminate the sample. Patients should be told to avoid zinc supplements before the test only if the clinical question concerns baseline status, but recent supplementation is important to disclose. Because inflammation redistributes zinc out of the circulation, the presence of concurrent infection or acute illness should be noted for correct interpretation.

\section*{Contraindications and Precautions}
There are no absolute contraindications to standard venipuncture. Critical pre-analytical and interpretive cautions: (1) Contamination is the most common cause of spurious elevation --- avoid glass, rubber-stopper tubes, and skin contact with the sample; royal blue-top trace-element tubes are required. (2) Hemolysis invalidates the result because red blood cells contain roughly 10-fold higher zinc than serum. (3) Serum zinc falls 10--20\% during acute illness, surgery, or infection because zinc is redistributed to the liver and the acute-phase response; a low zinc in an acutely ill patient does not necessarily indicate true deficiency. (4) Recent meals lower serum zinc, so fasting samples are preferred. (5) Pregnancy, oral contraceptive use, and estrogen therapy lower serum zinc. (6) Diurnal variation: levels are highest in the morning. (7) Low albumin lowers total serum zinc simply because albumin is the major zinc-binding protein --- in hypoalbuminemic patients the result should be interpreted cautiously, and plasma zinc may be preferable. (8) Antiepileptic drugs, diuretics, penicillamine, and high-dose iron supplements can alter zinc levels.

\section*{Risks}
Minimal risk. Slight pain or bruising at the venipuncture site. Rarely: hematoma, infection, or vasovagal syncope.

\section*{What Happens After the Test}
Results are typically available within 3--7 days because trace-element testing is often performed at a reference laboratory. The result is interpreted alongside clinical findings, dietary history, serum albumin, and inflammatory markers. Low zinc is generally treated with oral zinc supplementation (typically 25--50 mg elemental zinc daily for adults, adjusted by weight in children) with repeat measurement after 4--12 weeks to confirm adequacy. High zinc levels warrant discontinuation of supplements and evaluation for occupational or accidental exposure, because chronic excess causes copper deficiency with anemia and neuropathy. Patients found deficient should have a follow-up level checked after the treatment course to guide maintenance dosing.

\section*{Understanding Results}

\subsection*{Below Normal}
\begin{itemize}
  \item \textbf{Zinc deficiency (serum zinc below approximately 70 mcg/dL in adults, with laboratory-specific cutoffs):} The most common causes are inadequate intake (restrictive diets, anorexia nervosa, parenteral nutrition without zinc), malabsorption (Crohn disease, celiac disease, short bowel syndrome, cystic fibrosis, chronic pancreatitis), and increased losses (chronic diarrhea, fistulas, burns, hemodialysis).
  \item \textbf{Clinical features of deficiency:} Impaired growth and hypogonadism in children and adolescents; alopecia; acrodermatitis-like skin rash around the mouth and extremities; impaired wound healing; increased susceptibility to infection; hypogeusia (impaired taste); night blindness; and in severe deficiency, acrodermatitis enteropathica.
  \item \textbf{Caution:} Serum zinc is falsely lowered by inflammation, acute illness, surgery, pregnancy, and hypoalbuminemia, so a low value must be interpreted in clinical context and may warrant confirmation with erythrocyte zinc or a repeat measurement after the acute illness resolves.
\end{itemize}

\subsection*{Above Normal}
\begin{itemize}
  \item \textbf{Serum zinc above the reference range is usually due to contamination at collection, recent oral supplementation, or hemolysis of the specimen.}
  \item \textbf{Chronic excessive zinc supplementation:} Causes symptomatic copper deficiency (microcytic anemia, neutropenia, myelopathy with ataxia and spasticity) because zinc induces intestinal metallothionein, which binds copper and blocks its absorption.
  \item \textbf{Occupational or industrial exposure:} Inhalation of zinc oxide fumes causes metal fume fever; chronic exposure is evaluated by urine zinc.
  \item \textbf{Acrodermatitis enteropathica on excessive replacement therapy, and rare inborn errors of zinc metabolism.}
  \item \textbf{Associated with hemolysis of the sample and with extended tourniquet time.}
\end{itemize}

\section*{Normal Results}
\begin{itemize}
  \item \textbf{Serum zinc (adults):} 70--120 mcg/dL (10.7--18.4 $\mu$mol/L), depending on the laboratory
  \item \textbf{Plasma zinc:} Similar to serum but slightly higher; plasma is preferred in some reference laboratories to reduce platelet-release contamination
  \item \textbf{Erythrocyte zinc:} Approximately 12--15 $\mu$g/mL red cells; reflects longer-term status
  \item \textbf{Urine zinc (24-hour):} 150--1,200 mcg/24 hours (varies with intake)
  \item \textbf{Conversion:} 1 $\mu$mol/L zinc = 6.54 mcg/dL
\end{itemize}

\section*{Follow-up Tests}
\begin{itemize}
  \item \textbf{Serum copper and ceruloplasmin:} Zinc and copper are physiologically linked; measure together in malabsorption or when toxicity is suspected, because excess zinc causes copper deficiency
  \item \textbf{Serum albumin:} Required for interpretation of total serum zinc because albumin carries roughly 70\% of circulating zinc
  \item \textbf{Inflammatory markers (CRP or ESR):} A low zinc in the setting of acute-phase inflammation may be spurious; documenting inflammation guides interpretation
  \item \textbf{Plasma or erythrocyte zinc:} Alternative measures when serum zinc is ambiguous or contamination is suspected
  \item \textbf{24-hour urine zinc:} Quantifies gastrointestinal losses (Crohn disease, short bowel syndrome) and monitors treatment of acrodermatitis enteropathica
  \item \textbf{Other trace elements (selenium, magnesium):} Deficiencies commonly coexist in malabsorption and parenteral nutrition
  \item \textbf{Complete blood count:} Screening for the anemia and neutropenia of copper deficiency if chronic zinc excess is suspected
  \item \textbf{Repeat serum zinc after 4--12 weeks of supplementation:} Confirms repletion and guides maintenance dosing
\end{itemize}

\section*{References}
\begin{enumerate}
  \item Prasad AS. Discovery of human zinc deficiency: its impact on human health and disease. Adv Nutr. 2013;4(2):176--190.
  \item Wessells KR, Brown KH. Estimating the global prevalence of zinc deficiency. PLoS One. 2012;7(11):e50568.
  \item Gibson RS, Hess SY, Hotz C, Brown KH. Indicators of zinc status at the population level. Food Nutr Bull. 2008;29(2 Suppl):S154--S163.
  \item Hambidge KM. Zinc and health: current status and future directions. J Nutr. 2000;130(5S Suppl):1344S--1349S.
  \item Krebs NF. Update on zinc deficiency and excess in clinical pediatric practice. Ann Nutr Metab. 2013;62(Suppl 1):19--29.
\end{enumerate}

\clearpage
\section*{Regulatory Status and Authorization Date}
Regulatory status applies to the specific assay, instrument, manufacturer, and intended use rather than to the test name alone. No single approval date can be assigned to this test category. For a commercial assay, verify the current product in the relevant regulator's database: the U.S. Food and Drug Administration (FDA) clearance, approval, or authorization; India's Central Drugs Standard Control Organisation (CDSCO) licence or permission; the European Union In Vitro Diagnostic Regulation (EU IVDR) CE certificate issued through the applicable conformity-assessment route; or the United Kingdom Medicines and Healthcare products Regulatory Agency (MHRA) registration and UKCA/CE status. Laboratory-developed tests may not have an individual FDA, CDSCO, EU IVDR, or MHRA approval date.
\clearpage
